
\subsection{Grupių prezentacijos}

Grupės prazentacija yra būdas apibrėžti grupę,
pateikiant jos generatorius ir sąryšius kuriuos jie tenkina.
Formaliai tokia grupė yra faktor-grupė laisvosios grupės,
generuojamos duotų generatorių.
Toliau pristatysime kobinatorinės grupių teorijos pagrindus,
reikalingus apibrėžti grupės prezentacija.

Tegul $S = \qty{a,b,c,\cdots}$ yra aibė generatorių.
Literatūroje įprasta aibę $S$ vadinti \emph{abėcėle}, 
o jos elementus - \emph{raidėmis}.
Tegul $S^{-1} = \qty{a^{-1},b^{-1},c^{-1},\cdots}$
yra aibė sudaryta iš \emph{atvikštinių raidžių}
(kol kas šios raidės yra tik aibės $S^{-1}$ elementai be papildomos struktūros).
Aibę $S \cup S^{-1}$ vadinsime \emph{išplėstine abėcėle}.
Baigtines sekas $ f_1 f_2 \cdots f_n$ vadinsime \emph{žodžiais},
čia visi $f_i \in S \cup S^{-1}$.
Tegul $e$ žymi tuščia žodį,
tai yra žodį sudarytą iš nulio raidžių.
Tegul tarp dviejų žodžių yra ekvivalntiškumo sąryšis
$f_1 f_2 \cdots f_n \sim f'_1 f'_2 \cdots f'_m$,
jei iš vieno žodžio galime gauti kitą
pakartotinai įterpdami raidės ir jos anvikštinės raidė poras $xx^{-1}$ ir $x^{-1}x$, $x \in S$.
Visi žodžiai tarp kurių yra ekvivalentiškumo sąryšis sudaro ekivalentiškumo klasę.
\begin{definition}
	Laisvoji grupė $F_S$ generuojama \emph{abėcėlės} $S$ yra sudaryta iš
	\emph{žodžių} ekvivalentiškumo klasių,
	kartu su daugyba duodama \emph{žodžių} (klasės atstovų) sujungimo.
\end{definition}

\begin{example}
	Parodykime kad $F_S$ tenkina grupės aksiomas.
	\begin{itemize}
		\item Uždarumas:
			sujungus du žodžius gausime nauja baigtinę raidžių seką, kuri taip pat yra žodis.
		\item Asociatyvumas:
			sekų jungimas yra asociatyvus.
		\item Egzistuoja vienetinis elementas:
			betkurį žodžį sujungus su tuščiu žodžiu $e$, gausime nepakeistą žodį.
		\item Egzistuoja atvikštiniai elementai:
			žodžiui $f_1 f_2 \cdots f_n$ atvikštinis yra $f_{n}^{-1} f_{n-1}^{-1} \cdots f_1^{-1}$,
			kur 
			\begin{equation}
				f_i^{-1}= 
				\begin{cases}
					x^{-1}\,,\;	&\text{kai } f_i=x
					\\
					x\,,		&\text{kai } f_i=x^{-1}
				\end{cases}
				\,,\qquad
				x \in S
				\,.
			\end{equation}
			Šių žodžių sandauga priklauso tai pačiai ekvivalentiškumo klasei,
			kaip tuščias žodis $e$.
			\begin{equation}
				f_1 f_2 \cdots f_n f_{n}^{-1} f_{n-1}^{-1} \cdots f_1^{-1}
				\sim
				f_1 f_2 \cdots f_{n-1} f_{n-1}^{-1} f_{n-2}^{-1} \cdots f_1^{-1}
				\sim \cdots \sim 
				f_1 f_1^{-1} \sim e
			\end{equation}
	\end{itemize}
\end{example}

Laisvoji grupė $F_S$ yra bendriausia įmanoma grupė turinti elementus $\qty{a, b, c, \cdots}$.
Ją sudarėme visų žodžių aibei suteikdami sąryšio $xx^{-1}=e$ stuktūra ir asociatyvia uždara daugybą,
tai yra minimali struktūra reikalinga kad $F_S$ būtų grupė.
Toliau parodysime, kaip galima sudaryti šios grupės faktor-grupes,
kurios turi papildomą struktūrą.

Yra aibė $R \subset F_S$,
jos elementai vadinami \emph{saryšio žodžiais} (ang. \emph{relators}).
Tai bus žodžiai, kurie faktor-grupėje bus sutapatinti su vienetiniu elementu.
Tegul $N(R)$ yra aibės $R$ normalusis uždarinys (ang. \emph{normal closure})
\begin{equation}
	N(R) := \qty{
		\left.
		\prod_{i=1}^n g_i r_i g_i^{-1}
		\;\right|\;
		n \in \set N ,\,
		r_i \in R \cup R^{-1} ,\,
		g_i \in F_S
	}
	\,.
\end{equation}

\begin{proposition}
	$N(R)$ yra grupės $F_S$ normalusis pogrupis.
\end{proposition}
\begin{proof}
	Parodome pogrupio uždarumą.
	$N(R)$ yra visų sąryšio žodžių jungtinių elementų baigtinės sandaugos.
	Sudauginus dvi tokias sandaugas gaunama nauja baigtinė sandauga,
	tad $N(R)$ yra uždara daugybos atžvilgiu.

	Parodome, kad $gN(R)g^{-1} \subseteq N(R)$ su visais $g \in F_S$.
	\begin{equation}
		g\qty(\prod_{i=1}^n g_i r_i g_i^{-1})g^{-1} 
		= \prod_{i=1}^n g g_i r_i g_i^{-1}g^{-1}
		= \prod_{i=1}^n g g_i r_i (g g_i)^{-1}
		\,,
	\end{equation}
	elementai $g g_i$ priklauso grupei $F_S$,
	tad pagal $N(R)$ apibrėžimą $g\qty(\prod_{i=1}^n g_i r_i g_i^{-1})g^{-1} \in N(R)$.
\end{proof}


\begin{definition}
	Turint \emph{abėcėlę} $S$ ir \emph{sąryšio žodžius} $R$,
	grupės prezentacija $\present{S}{R}$ yra faktor-grupė $F_S/N(R)$.
\end{definition}

\begin{remark}
	Literatūroje yra dažniau naudojamas žymėjimas $\langle\,S\,|\,R\,\rangle$,
	bet mes pasirinkome kitokį žymėjimą,
	kad nekiltų rizikos sumaišyti grupės prezentacija su Dirac'o ,,bra-ket'' žymėjimu.
	Taip pat abėcėlė $S$ bus pateikta nekaip aibė, o elementais.
	Sąryšio žodžiai $R$ bus pateikti kaip grupę apibrėžiantys sąryšiai
	$r' = r$, kur $r'r^{-1}\in R$.
	Sąryšiai yra patogūs, nes $F_s$ elemetai $r$ ir $r'$ priklauso
	tai pačiai gretutinei klasei
	\begin{equation}
		r\, N(R)
		= r\, r^{-1}	(r'r^{-1})	r \,N(R)
		= r'\,N(R)
		\,,
	\end{equation}
	tad grupėje $\present{S}{\qty{r'r^{-1},\cdots}}$
	turime tapatybę $r' = r$ (rašome tik gretutinės klasės atstovus).
\end{remark}

\begin{example}
	Grupės prezentaciją
	$ \present{S}{R} $, $ S = \qty{a, b} $,
	$ R = \qty{a^3, b^2, aba^{-1}b^{-1}} $,
	rašysime, taip 
	\begin{equation}
		\present{a,b}{a^3 = b^2 = e ,\, ab = ba}
	\end{equation}
	ir su grupės elelementais galime atlikti skaičiavimus kaip kad
	\begin{equation}
		abab = abba = aea = a^2
		\,.
	\end{equation}
\end{example}


\begin{remark}
	Svarbu paminėti, kad grupės presentacija nėra unikali.
	\begin{itemize}
		\item Galime pridėti sąryšio žodį $w \in N(R)$, $w \notin R$
			ir gausime lygias grupių prezentacijas
			\begin{equation}
				\present S R = \present{S}{R \cup \qty{w}}
				\,.
			\end{equation}
		\item Jei buvo pasirinkta daugiau sąryšio žodžių negu būtina,
			galima dalį sąryšio žodžių išmesti iš $R$.
		\item Galime pridėti generatorių $d$ prie abėcėlės $S$
			ir naują sąryšio žodį $dw^{-1}$ prie $R$,
			kur $w$ yra bet koks žodis iš piminės abėcėlės $S$ raidžių,
			ir gausime izomorfiškas grupes
			\begin{equation}
				\present S R \cong \present{ S \cup \qty{d} }{ R \cup \qty{dw^{-1}} }
				\,.
			\end{equation}
		\item Jei buvo pasirinkta daugiau generatorių negu būtina,
			galima dalį generatorių pakeisti žodžiais,
			sudarytais iš likusių generatorių.
	\end{itemize}
	Šie grupių prezentacijų pakeitimai vadinami Tieze transformacijomis.
	Klausimai, ar dvi duotos grupių prezentacijos yra izomorfiškos,
	ir ar duota grupės prezentacija yra minimali
	(neturi generatorių ar sąryšio žodžių,
	kuriuos išmetus gautume izomorfišką grupę), -
	yra netrivialūs kombinatorinės grupių teorijos uždaviniai.
\end{remark}

\begin{example}
	Parodysime, kad grupės prezentacija $\present{a,b}{a^2=b^2=ab=e}$
	yra izomorfiška trivialiai grupei $E$.
	\begin{equation}
		\begin{aligned}
			ab=e	&\;\;\text{ir}\;\; a^2=e	&\implies&& a &= b		\,,\\
			a=b		&\;\;\text{ir}\;\; b^3=e	&\implies&& a^3 &= e	\,,\\
			a^3=e	&\;\;\text{ir}\;\; a^2=e	&\implies&& a &= e		\,,\\
			a=e		&\;\;\text{ir}\;\; a=b		&\implies&& b &= e		\,.
		\end{aligned}
	\end{equation}
	Abu generatoriai yra lygūs vienetiniam elementui,
	tad $ \present{a,b}{a^2=b^2=ab=e} \cong E $.
\end{example}


